\section{由算符乘积展开得到的因子化}
\label{sec:ope}

本节我们利用算符乘积展开推导准 PDF 的匹配关系，以及\spcorr 和赝 PDF 的等价匹配关系。为简单起见，这三个等价情形分别在单独的小节中给出。

\subsection{空间关联函数的 OPE 与因子化}

OPE 是一种在欧氏极限 $z^2\to 0$ 下把间隔为 $z^\mu$ 的非定域算符展开为定域算符的技术。
它既适用于裸的正则化算符，也适用于重正化算符；我们的焦点是后者。对规范不变的 Wilson 算符 $\tilde{O}_\Gamma(z)$，已经证明它在坐标空间可以作乘法重正化~\cite{Ji:2017oey,Ishikawa:2017faj}
\beq \label{eq:ren}
\tilde{O}_\Gamma(z,\mu) = Z_{\psi,z}\, e^{\delta m |z|} \tilde{O}_\Gamma(z,\epsilon)\,,
\eeq
其中 $\delta m$ 概括了 Wilson 线自能带来的幂次发散，$Z_{\psi,z}$ 只依赖端点 $z,0$ 并重正化对数发散。这一乘法重正化在更早的文献~\cite{Ishikawa:2016znu,Chen:2016fxx,Constantinou:2017sej}中也有讨论。为简单起见，本节对式~(\ref{eq:ren})取 $\Gamma=\gamma^z$。

在 $\overline{\rm MS}$ 方案中，幂次发散消失，利用 OPE，当 $z^2\to 0$ 时重正化的 $\tilde{O}_\Gamma(z,\mu)$ 可以按定域规范不变算符展开，得到
\begin{align}\label{eq:ope-tq}
\tilde{O}_{\gamma^z}(z,\mu) = & \sum_{n=0}^\infty \left[C_n ({\mu}^2 z^2)\frac{(-iz)^n}{n!} e_{\mu_1}\cdots 				e_{\mu_n}O_1^{\mu_0\mu_1\cdots\mu_n}(\mu)\right.
\\
&+C'_n ({\mu}^2 z^2)\frac{(-iz)^n}{n!} e_{\mu_1}\cdots e_{\mu_n}O_2^{\mu_0\mu_1\cdots\mu_n}(\mu)\nn\\
& +  \text{高扭度算符}\Big]\,,\nn
\end{align}
其中 $\mu_0=z$，$C_n=1+O(\alpha_s)$ 与 $C'_n=O(\alpha_s)$ 是 Wilson 系数，而 $O_1^{\mu_0\mu_1\cdots\mu_n}(\mu)$ 和 $O_2^{\mu_0\mu_1\cdots\mu_n}(\mu)$ 是 OPE 中领头幂次仅有的容许的、重正化的无迹对称扭度-2 夸克与胶子算符，
\begin{align}\label{eq:twist-2}
O_1^{{\mu}_0 {\mu}_1 \ldots {\mu}_n}(\mu) = & Z_{n+1}^{qq} \bigl(
 \bar{\psi} \gamma^{({\mu}_0 } iD^{{\mu}_1} \cdots iD^{ {\mu}_n)} \psi -
\text{迹项} \bigr) \,,
  \\
O_2^{{\mu}_0 {\mu}_1 \ldots {\mu}_n}(\mu) = & Z_{n+1}^{qg}\bigl(
 F^{(\mu_0\rho} iD^{{\mu}_1} \cdots iD^{ {\mu}_{n-1}} F_\rho^{~\mu_n)} - \text{迹项} \bigr)
  .\nn
\end{align}
这里 $Z_{n+1}^{ij}=Z_{n+1}^{ij}(\mu,\epsilon)$ 是乘法 $\overline{\rm MS}$ 重正化因子，$(\mu_0 \cdots \mu_n)$ 表示这些洛伦兹指标的对称化。

上述 OPE 对算符本身成立，其中我们隐含地把自己限制在扭度展开适用的矩阵元子空间。在同位旋矢量情形，与胶子算符的混合不存在，本文其余部分均限于该情形。当 $O_1^{\mu_0\mu_1\cdots\mu_n}$ 在核子态中取矩阵元时，
\beq \label{eq:def-moments}
\langle P | O_1^{{\mu}_0 {\mu}_1 \cdots {\mu}_n} | P\rangle =  2a_{n + 1}(\mu)\left(P^{{\mu}_0} P^{{\mu}_1} \ldots P^{{\mu}_n} - \text{迹项}\right), 
\eeq
其中 $a_{n+1}(\mu)$ 是 PDF 的 $(n+1)$ 阶矩，
\beq \label{eq:moment}
a_{n + 1} \left(\mu\right)=\int_{-1}^1 dx\,x^n q \left(x,\mu\right)\,,
\eeq
而迹项的显式表达式已在文献~\cite{Nachtmann:1973mr,Georgi:1976ve,Chen:2016utp}中导出。与式~(\ref{eq:moment})的逆关系是：$q(x,\mu)$ 具有正比于 $\delta$ 函数 $n$ 阶导数的项（如 $\delta^{(n)}(x)\, a_n(\mu)$）的展开，不含任何非平凡的短距离 Wilson 系数。

正如文献~\cite{Ji:2017rah}所指出的，要在小 $z^2$ 下、在 $|\zeta|= |P^z z|\sim 1$ 处为\spcorr 获得足够的信息，我们必须选择远大于标度 $\Lambda_{\rm QCD}$ 的 $P^z$。当 $P_z^2\gg\{\Lambda_{\rm QCD}^2, M^2\}$ 时，式~(\ref{eq:def-moments})中的迹项被 $M^2/P_z^2$ 的幂压低，而式~(\ref{eq:ope-tq})中高扭度算符的贡献被 $\Lambda_{\rm QCD}^2/P_z^2$ 或 $z^2\Lambda_{\rm QCD}^2$ 的幂压低。因此，扭度-2 贡献是大动量下核子矩阵元 $\langle P|\tilde{O}_{\gamma^z}(z)|P\rangle$ 的领头近似。从现在起我们在讨论中略去所有幂次修正。


$\tilde{O}_{\gamma^z}(z)$ 的 OPE 中的 Wilson 系数 $C_n ({\mu}^2 z^2)$ 可以在 $\mu\sim 1/|z|\gg \Lambda_{\rm QCD}$ 的微扰论中计算。在
$\overline{\ensuremath{\operatorname{MS}}}$ 方案中，$C_n$ 在 $z^2 =
0$ 附近呈对数奇异性，$\langle P|\tilde{O}_{\gamma^z}(z,\mu)|P\rangle$ 也是如此。正因如此，准 PDF $\tilde q(x,P^z,\mu)$ 的 $x$ 阶矩正比于 $C_n|_{z=0}$——后者是发散的——准 PDF 在无穷 $P^z$ 极限下并不会简单地变成 PDF。相反，我们需要一个把准 PDF 匹配到 PDF 的因子化公式。对比之下，赝 PDF 的矩确实存在，因为在取 $x$ 阶矩时我们保持 $\mu^2 z^2$ 固定。然而我们仍需要一个把赝 PDF 匹配到 PDF 的因子化公式。下面我们对此作进一步评论。

基于式~(\ref{eq:ope-tq}-\ref{eq:moment})，可以把\spcorr 的领头扭度近似写为
\begin{align} \label{eq:fac-tQ-orig}
\tilde{Q}_{\gamma^z} \bigl(\zeta, {\mu}^2 z^2\bigr)
=&  \sum_{n} C_n ({\mu}^2 z^2)  \frac{(-i \zeta)^n}{n!} a_{n + 1}
 \left(\mu\right)
 \\
= & \sum_{n} C_n (\mu^2z^2)
\frac{(-i \zeta)^n}{n!}  \int_{-1}^1 dy\,y^n q \left( y, \mu\right)
 \,.  \nn
\end{align}
应当指出，到目前为止我们所作的唯一近似是忽略被小 $z^2$ 以及核子大动量 $P^z$ 压低的高扭度效应。在 $P_z^2\gg M^2,\Lambda_{\rm QCD}^2$ 的极限下有 $P^0\sim P^z$，因此即使式~(\ref{eq:ope-tq})中 $\mu_0=0$，$\tilde{O}_{\gamma^0}(z)$ 的领头近似仍由式~(\ref{eq:fac-tQ-orig})中的扭度-2 贡献给出，只是系数 $C_n$ 有所修改。

基于 OPE 结果式~(\ref{eq:fac-tQ-orig})，我们可以为欧氏\spcorr 导出一个因子化公式。首先定义函数 $\mathcal{C} (\alpha, {\mu}^2 z^2)$：
\begin{align}  \label{eq:fac-C}
\mathcal{C} (\alpha, {\mu}^2 z^2) \equiv \int\! \frac{d \zeta}{2 \pi} \,
 e^{i \alpha \zeta}  \sum_n C_n ({\mu}^2 z^2)  \frac{(-i \zeta)^n}{n!}
 \,.
\end{align}
由式~(\ref{eq:fac-tQ-orig})以及式~(\ref{eq:pseudo})中傅里叶变换关系的重正化类比可知，$\mathcal{C}(\alpha,\mu^2z^2)$ 对应于 $a_{n+1}(\mu)=1$ 这一特例下的赝 PDF。文献~\cite{Radyushkin:1983wh,Radyushkin:2016hsy}的分析蕴含 $\mathcal{C}(\alpha,\mu^2z^2)$ 的支撑区为 $-1 \le \alpha \le 1$。注意到
\begin{align}
  \int d\alpha\, e^{-i\alpha(y \zeta)}\, \mathcal{C} (\alpha, {\mu}^2 z^2)
   &= \sum_n C_n(\mu^2 z^2) \frac{(-i\zeta y)^n}{n!} \,,
\end{align}
我们从式~(\ref{eq:fac-tQ-orig})得到
\begin{align}
\tilde{Q}(\zeta,\mu^2 z^2)
= &\int_{-1}^1\!\! dy\int_{-1}^1\!\! d\alpha\ e^{-i\alpha(y\zeta)}\, \mathcal{C}(\alpha,\mu^2z^2)\, q(y,\mu) \,.
\end{align}
最后，利用式~(\ref{eq:Qdefn})的重正化类比的逆变换，
\begin{align} \label{eq:FTq}
Q(\zeta,\mu) &= \int_{-1}^1\!\! dy\: e^{-i y\zeta}\,
  q(y,\mu) \,.
\end{align}
我们得到
\begin{align} \label{eq:io-Q-fact}
\tilde{Q}(\zeta,\mu^2 z^2)
=& \int_{-1}^1 d\alpha\ \mathcal{C}(\alpha,\mu^2z^2)\, Q(\alpha \zeta,\mu)
 + {\cal O}(z^2\Lambda_{\rm QCD}^2)
\,.
\end{align}
式~(\ref{eq:io-Q-fact})就是可在格点上计算的\spcorr $\tilde Q(\zeta,\mu^2 z^2)$ 的因子化公式，它把它表示成定义 PDF 的光锥关联 $Q(\zeta,\mu)$。它与文献~\cite{Braun:2007wv,Bali:2017gfr}中用于计算 $\pi$ 介子分布振幅的\spcorr 因子化公式具有相同的结构。

\subsection{准 PDF 的因子化}

重正化准 PDF 定义为重正化\spcorr 的傅里叶变换，
\begin{align} \label{eq:fac-tq-orig0}
\tilde{q} \Bigl(  x, \frac{\mu}{P_z}\Bigr)
\equiv &\int \frac{d \zeta}{2 \pi}\: e^{ix \zeta}\:  \tilde{Q} \biggl(\zeta, \frac{{\mu}^2\zeta^2}{P_z^2}\biggr) \,.
\end{align}
注意这里既可以使用 $\tilde{Q}_{\gamma^z}$ 也可以使用 $\tilde{Q}_{\gamma^0}$。
利用式~(\ref{eq:fac-tQ-orig})的结果，这给出
\begin{align} \label{eq:fac-tq-orig}
\tilde{q} \Bigl( & x, \frac{\mu}{P_z}\Bigr)
 \\
= &  \int_{-1}^1 dy \biggl[ \int \frac{d \zeta}{2 \pi} e^{ix \zeta}
\sum_{n=0} C_n \Bigl( \frac{{\mu}^2\zeta^2}{P_z^2}\Bigr)  \frac{(-i
	\zeta)^n}{n!} y^n \biggr] q \left( y, \mu \right)
\nn\\
= &  \int_{-1}^1 {dy\over |y|} \biggl[\int \frac{d \zeta}{2 \pi} e^{i
	\frac{x}{y} \zeta}  \sum_{n=0} C_n \Bigl( \frac{{\mu}^2\zeta^2 }{(yP^z)^2}\Bigr)\,  \frac{(-i \zeta)^n}{n!}  \biggr] q \left( y,\mu \right) \,.
  \nn
\end{align}
已经可以看出，匹配核是 $x / y$
和 $\mu / (|y|P^z)$ 的函数。我们把该核定义为
\begin{align}   \label{eq:fac-c}
C \left( {x\over y}, \frac{\mu}{|y|P^z} \right)
   \equiv \int\! \frac{d \zeta}{2 \pi} \: e^{i {x\over y} \zeta}
   \sum_{n=0} C_n \Bigl(
\frac{{\mu}^2\zeta^2}{(yP^z)^2}\Bigr) \, \frac{(-i \zeta)^n}{n!}
 \,,
\end{align}
于是式~(\ref{eq:fac-tq-orig})可以改写为
\begin{align} \label{eq:fac-tq}
\tilde{q} \left( x, \frac{\mu}{P^z}\right)
  = \int_{-1}^1 \frac{dy}{|y|}\: C \Bigl(
   \frac{x}{y}, \frac{\mu}{|y|P^z} \Bigr)\: q \left( y,\mu\right)
   \,,
\end{align}
这就是准 PDF 的 $\overline{\rm MS}$ 因子化公式。这一结果表明，式~(\ref{eq:momfact})中的因子化公式在 $\overline{\rm MS}$ 方案下必须取 $p^z=|y|P^z$。我们将在\sec{ren}证明这对任何准 PDF 重正化方案都成立。这与早期准 PDF 论文中猜想并使用的取法 $p^z=P^z$ 不同~\cite{Ji:2013dva,Ji:2014gla,Xiong:2013bka}。物理上，式~(\ref{eq:fac-tq})中的正确结果可以这样理解：匹配系数只对微扰的部分子动力学敏感，因此出现的是部分子动量的大小 $|y| P^z$ 而不是强子动量 $P^z$。

利用式~(\ref{eq:fac-tq-orig0})取准 PDF 的矩得
\begin{align} \label{eq:momentqpdf}
 &\int_0^1\!\! dx\: x^n \tilde{q} \left( x, \frac{\mu}{P^z}\right)
= \Big(i\frac{d}{d\zeta}\Big)^n \tilde{Q} \biggl(\zeta, \frac{{\mu}^2\zeta^2}{P_z^2}\biggr) \Bigg|_{\zeta\to 0}
 \\
& \quad =  \sum_{n'}  \Big(i\frac{d}{d\zeta}\Big)^n \bigg[ C_{n'} \Big(\frac{{\mu}^2 \zeta^2}{P_z^2}\Big)  \frac{(-i \zeta)^{n'}}{n'!} \bigg] \Bigg|_{\zeta\to 0}  a_{n' + 1}
\left(\mu\right) \,.  \nn
\end{align}
由于 $C_{n'}$ 系数具有 $\ln(\zeta^2)$ 依赖性，$n'=n$ 的导数在 $\zeta\to 0$ 时总有对数奇异性，而 $n'<n$ 时还会有更奇异的项。这解释了为什么短距离 Wilson 系数导致准 PDF 的矩不存在。


\subsection{赝 PDF 的因子化}

重正化赝 PDF 是重正化\spcorr 的傅里叶变换：
\begin{align} \label{eq:pseudoRen}
\mathcal{P} \left( x, \mu^2 z^2\right)
   = \int_{-\infty}^\infty \! \frac{d \zeta}{2\pi}\:
    e^{ix \zeta}\: \tilde{Q}\bigl( \zeta, \mu^2 z^2 \bigr)
   \,.
\end{align}
由于赝 PDF 和\spcorr 都以与 $\zeta$ 无关的方式乘法重正化，这可以由式~(\ref{eq:pseudo})立即得到。
若对式~(\ref{eq:io-Q-fact})把\spcorr $\tilde Q(\zeta,\mu^2z^2)$ 傅里叶变换成赝 PDF，把光锥关联 $Q(\alpha\zeta,\mu)$ 变换成 PDF，则立即可得赝 PDF 的因子化公式，
\begin{align}\label{eq:ps-q-fact}
\mathcal{P}(x,z^2\mu^2) =
 &\int_{|x|}^1 {dy\over |y|}\
  \mathcal{C}\left({x\over y},\mu^2 z^2\right) q(y,\mu)\\
&+\int^{-|x|}_{-1} {dy\over |y|}\
  \mathcal{C}\left( \frac{x}{y}, \mu^2z^2\right) q(y,\mu)
  \nn \\
 &+ {\cal O}(z^2\Lambda_{\rm QCD}^2)\,,\nn
\end{align}
这正是式~(\ref{eq:pseudofact})的小 $z^2$ 因子化公式。式~(\ref{eq:ps-q-fact})中积分的上下限直接来自匹配系数 $\mathcal{C}(\alpha,z^2\mu^2)$ 的支撑区 $-1\le \alpha \le 1$，并且提醒左边赝 PDF 也满足 $-1\le x\le 1$。

由于赝 PDF 的 $x$ 范围有界，式~(\ref{eq:pseudoRen})中指数的级数展开各项存在，
\begin{align}
  \tilde Q(\zeta,\mu^2 z^2)
   = \sum_{n=0}^\infty \int_{-1}^{1}\!\! dx\: \frac{(-i\zeta)^n x^n}{n!}
     {\cal P}(x,\mu^2 z^2) \,.
\end{align}
与式~(\ref{eq:fac-tQ-orig})比较可知，赝 PDF 的矩为
\begin{align}  \label{eq:Pmoment}
  \int_{-1}^{1}\!\! dx\: x^n \:  {\cal P}(x,\mu^2 z^2)
  &= C_n(\mu^2 z^2)\, a_{n+1}(\mu) \,.
\end{align}

至此，我们证明了准 PDF 的大 $P^z$ 因子化以及空间关联函数和赝 PDF 的小 $z^2$ 因子化。推导出其中一个因子化之后，其余的便随即得到，因为它们只是同一空间关联函数的不同表示。的确我们看到，在领头幂次下准 PDF 与赝 PDF 通过其定义相联系：
\beq \label{eq:equiv}
\tilde{q}\Bigl(x,{\mu\over P^z}\Bigr)
  = \int\! {d\zeta\over 2\pi}\ e^{ix\zeta}\int_{-1}^1\!\! dy \ e^{-iy\zeta}
  \: \mathcal{P}\biggl(y,{\mu^2 \zeta^2\over P_z^2}\biggr)\,,
\eeq
其中用到 $z=\zeta/P^z$。
基于式~(\ref{eq:fac-C})与式~(\ref{eq:fac-c})，两个因子化定理中的 Wilson 系数也保持同样的关系：
\begin{align} \label{eq:ftc-C}
 C\left(\xi,{\mu\over |y|P^z}\right)
  &= \int\! {d\zeta\over 2\pi}\ e^{i\xi\zeta}\int_{-1}^1\!\! d\alpha\: e^{-i\alpha\zeta}\: \mathcal{C}\left(\alpha,{\mu^2 \zeta^2\over (yP^z)^2}\right) .
\end{align}
对于式~(\ref{eq:equiv})与(\ref{eq:ftc-C})中的关系，准 PDF 与赝 PDF（或其相应系数）应采用同样的 $\Gamma=\gamma^0$ 或 $\gamma^z$ 选择。

总之，存在唯一的因子化公式把准 PDF、\spcorr 和赝 PDF 匹配到 PDF。既然它们到 PDF 的因子化都需要小距离和大核子动量，它们格点计算的设定也必然相同。因此，LaMET 方法与赝分布方法原则上是彼此等价的，它们也许只在格点实现相关的效应上有所差别。

文献~\cite{Radyushkin:2017cyf}中猜测，可以在格距为 $a$ 的格点上研究比值函数
\beq \label{eq:ratio}
\tilde{Q}(\zeta,z^2,a^{-1})/\tilde{Q}(0,z^2,a^{-1})
\eeq
且 $O(z^2)$ 修正可能近似抵消。这一想法在文献~\cite{Orginos:2017kos}的格点 QCD 中得到检验，结果表明在小值处该比值随 $z^2$ 缓慢演化。那么从这个比值中能够提取哪类非微扰信息就是一个有趣的问题。

这个问题可以用\spcorr 的小 $z^2$ 因子化来回答。
根据式~(\ref{eq:fac-tQ-orig})，
\beq
	\tilde{Q}(0,\mu^2 z^2) = C_0(\mu^2z^2) + O(z^2\Lambda_{\rm QCD}^2)\,,
\eeq
其中在 $\overline{\rm MS}$ 方案下单圈为
\beq  \label{eq:C0}
	C_0(\mu^2z^2) = 1 + {\alpha_sC_F\over 2\pi} \biggl[ {3\over2} \ln(\mu^2z^2 e^{2\gamma_E}/4)+{5\over2} \biggr] \,,
\eeq
这最近也在文献~\cite{Zhang:2018ggy}中被导出。
于是比值变为
\begin{align} \label{eq:Qtratio}
	\frac{\tilde{Q}\left(\zeta, {\mu}^2 z^2\right)}{\tilde{Q}\left(0, {\mu}^2 z^2\right)}
	=&\sum_n \frac{C_n (\mu^2z^2)}{C_0 (\mu^2z^2)}	\frac{(-i \zeta)^n}{n!} a_{n+1}(\mu)\\
	&+ O(z^2\Lambda_{\rm QCD}^2)\,.\nn
\end{align}
利用式~(\ref{eq:Pmoment})以及下文式~(\ref{eq:1looppseudopdfb})中我们的 $\overline{\rm MS}$ 单圈微扰赝 PDF 结果，对 $\Gamma=\gamma^0$ 可得
\begin{align}
  C_n(\mu^2 z^2) &= 1 +\frac{\alpha_s C_F}{2\pi} \biggl[ \Bigl(
  \frac{3\!+\!2n}{2\!+\!3n\!+\!n^2}
  \!+\! 2H_n\Bigr) \ln\frac{\mu^2z^2 e^{2\gamma_E}}{4}
  \nn\\
 &\quad\quad
  +   \frac{5\!+\!2n}{2\!+\!3n\!+\!n^2}+2(1-H_n)H_n - 2 H_n^{(2)}
  \biggr] ,
\end{align}
其中调和数为 $H_n=\sum_{i=1}^n 1/i$ 与 $H_n^{(2)} =\sum_{i=1}^n 1/i^2$。对 $\Gamma=\gamma^z$ 情形有 $C_n^{\gamma^z}(\mu^2 z^2) = C_n(\mu^2 z^2) + \Delta C_n^{\gamma^z}(\mu^2 z^2)$，其中
\begin{align}
  \Delta C_n^{\gamma^z}(\mu^2 z^2)
   &= \frac{\alpha_s C_F}{2\pi} \: \frac{2}{2\!+\!3n\!+\!n^2} \,,
\end{align}
这也会修改 $n=0$ 时的式~(\ref{eq:C0})。
在 $\mu\simeq 1/|z|$ 微扰展开适用的小 $z^2$ 区，式~(\ref{eq:Qtratio})中的比值对 $|z|$ 只有弱的对数依赖，这与文献~\cite{Orginos:2017kos,Karpie:2017bzm}的格点发现一致。对 $|z|$ 的弱依赖可以用 $\ln z^2$ 的演化方程定量描述~\cite{Radyushkin:2017cyf,Radyushkin:2017lvu}。根据我们的 OPE 分析，这里 $\tilde{Q}(0,\mu^2 z^2)$ 只充当被高扭度修正污染的整体归一化因子，而 $\ln z^2$ 演化可以准确地纳入式~(\ref{eq:io-Q-fact})的因子化公式之中，使我们能够从比值函数提取 PDF。同样的观点最近也在文献~\cite{Zhang:2018ggy}的工作中得到论证，该工作与本文同时出现。
